\documentclass[12pt,a4paper]{article}
\usepackage[utf8]{inputenc}
\usepackage[T1]{fontenc}
\usepackage{amsmath}
\usepackage{booktabs}
\usepackage{geometry}
\usepackage{hyperref}
\usepackage{cite}
\geometry{a4paper,top=2.5cm,bottom=2.5cm,left=2.5cm,right=2.5cm}
\hypersetup{colorlinks=true,linkcolor=blue,citecolor=blue,urlcolor=blue}

\title{\textbf{Self-Sustaining Tritium Regeneration Loop in\\
Proton--Electron Capture Chain Reactions:\\
A Closed Fuel Cycle for Fusion Energy}}
\author{Eymen Aydoğan}
\date{May 19, 2026}

\begin{document}
\maketitle

\begin{abstract}
We demonstrate theoretically that the proton--electron capture chain
mechanism (PECCNAL) produces a self-sustaining tritium regeneration
loop requiring no external tritium supply. Tritium consumed in
deuterium--tritium fusion (Step 3) is continuously regenerated via
two independent pathways: neutron--deuterium capture (Step 2) and
Lithium-6 neutron capture (Step 4). The net result is a closed fuel
cycle in which only deuterium and Lithium-6 are consumed, while
tritium acts as a catalytic intermediate. This resolves the central
fuel supply problem of D--T fusion energy.
\end{abstract}

\noindent\textbf{Keywords:} tritium breeding, closed fuel cycle,
self-sustaining loop, D--T fusion, neutron capture, lithium-6

\hrule\vspace{1em}

\section{Introduction}
Tritium ($^3$H) is an essential fuel for D--T fusion but is
radioactive (half-life 12.3 years) and not naturally abundant.
Current fusion programs rely on external tritium supply or blanket
breeding, neither of which achieves true fuel self-sufficiency.

In this paper, we show that the PECCNAL mechanism~\cite{aydogan2026a,aydogan2026b}
naturally produces a closed tritium loop, where tritium is
regenerated at the same rate it is consumed.

\section{The Tritium Loop}

The complete cycle is:
\begin{align}
p + e^- &\rightarrow n + \nu_e \quad -1.293\,\text{MeV} \tag{S1}\\
n + {}^2\text{H} &\rightarrow {}^3\text{H} + p \quad +2.220\,\text{MeV} \tag{S2}\\
{}^3\text{H} + {}^2\text{H} &\rightarrow {}^4\text{He} + n \quad +17.600\,\text{MeV} \tag{S3}\\
n + {}^6\text{Li} &\rightarrow {}^4\text{He} + {}^3\text{H} \quad +4.780\,\text{MeV} \tag{S4}
\end{align}

\subsection{Tritium Balance}

\begin{table}[h!]
\centering
\caption{Tritium production and consumption per cycle}
\vspace{0.5em}
\begin{tabular}{lcc}
\toprule
\textbf{Step} & \textbf{Tritium} & \textbf{Net} \\
\midrule
S2: $n + {}^2\text{H} \rightarrow {}^3\text{H}$ & +1 produced & \\
S3: ${}^3\text{H} + {}^2\text{H} \rightarrow {}^4\text{He}$ & -1 consumed & \\
S4: $n + {}^6\text{Li} \rightarrow {}^3\text{H}$ & +1 produced & \\
\midrule
\textbf{Net per cycle} & & \textbf{+1 surplus} \\
\bottomrule
\end{tabular}
\end{table}

The system produces \emph{more} tritium than it consumes,
achieving a tritium breeding ratio (TBR) $> 1$:

\begin{equation}
\text{TBR} = \frac{T_\text{produced}}{T_\text{consumed}}
= \frac{T_{S2} + T_{S4}}{T_{S3}} = \frac{1 + 1}{1} = 2
\end{equation}

A TBR of 2 means the system doubles its tritium inventory
per cycle, enabling exponential scale-up without external supply.

\section{Energy Balance of the Closed Loop}

\begin{equation}
E_\text{net} = -1.293 + 2.220 + 17.600 + 4.780
= \boxed{+23.307\,\text{MeV per cycle}}
\end{equation}

\section{Fuel Consumption}

Per complete cycle, the system consumes:
\begin{itemize}
\item 2 deuterium nuclei ($^2$H): one in S2, one in S3
\item 1 Lithium-6 nucleus in S4
\item 1 electron in S1 (replenished by ionization)
\end{itemize}

Tritium is \textbf{not consumed} on net — it is a catalytic
intermediate regenerated each cycle.

\section{Comparison with ITER Blanket Breeding}

\begin{table}[h!]
\centering
\caption{Tritium breeding comparison}
\vspace{0.5em}
\begin{tabular}{lcc}
\toprule
\textbf{Feature} & \textbf{ITER Blanket} & \textbf{PECCNAL Loop} \\
\midrule
TBR target & $>1.05$ & $2.0$ (theoretical) \\
External T supply needed & Yes (startup) & No \\
T regeneration pathway & Single (Li blanket) & Dual (S2 + S4) \\
Coulomb barrier & N/A & None \\
Scale-up & Linear & Exponential \\
\bottomrule
\end{tabular}
\end{table}

\section{Conclusion}
The PECCNAL mechanism achieves a self-sustaining tritium
regeneration loop with TBR = 2, requiring only deuterium and
Lithium-6 as consumed fuels. This eliminates the tritium supply
bottleneck that currently limits D--T fusion development.

\section*{Acknowledgments}
The author thanks the fusion physics community for published work
that informed this theoretical development.

\begin{thebibliography}{9}
\bibitem{aydogan2026a} E. Aydoğan, ``Coulomb-Barrier-Free Neutron Production via Proton--Electron Capture in Dense Hydrogen Gas with Lithium-6 Energy Extraction,'' arXiv preprint, 2026.
\bibitem{aydogan2026b} E. Aydoğan, ``Deuterium-Mediated Neutron Chain Fusion Initiated by Proton--Electron Capture with Lithium-6 Yield Enhancement,'' arXiv preprint, 2026.
\bibitem{pdg2022} R. L. Workman et al. (PDG), \textit{Prog. Theor. Exp. Phys.}, 083C01, 2022.
\end{thebibliography}
\end{document}
